Mapa průměrných mezd ve finančním sektoru a~pojišťovnictví (\ac{NACE} K) v~\acs{geo-EU27} umísťuje \acacc{ČR} do středního pole: Pražské finanční centrum vytváří nadprůměrné odměny ve srovnání s~celonárodním průměrem, avšak na mezinárodní úrovni zůstávají pod standardem Rakouska, Nizozemska nebo Lucemburska.
Finanční sektor je v~ČR charakteristický oligopolní strukturou se silným podílem zahraničních bank --- zaměstnavatelé v~tomto sektoru mají o~poznání vyšší ziskovost než v~jiných odvětvích, přesto \ac{KS} pokrytí je ve finančním sektoru velmi nízké.
Výsledkem je, že mzdová prémie finančního sektoru nad celostátním průměrem je v~\acloc{ČR} nižší než v~Rakousku nebo Švýcarsku, přestože produktivita na zaměstnance v~bankovnictví je komparabilní nebo vyšší --- přebytek jde primárně do dividendových výnosů zahraničních vlastníků.
Odvětvová \ac{KS} ve finančním sektoru by mohla zachytit část ziskového přebytku ve prospěch zaměstnanců --- nebyl by to administrativní zásah do trhu, ale institucionalizace kolektivní vyjednávací pozice zaměstnanců vůči oligopolně silným zaměstnavatelům.

Mapa průměrných mezd ve finančním sektoru a~pojišťovnictví (\ac{NACE} K) v~\acs{geo-EU27} umísťuje \acacc{ČR} do středního pole: Pražské finanční centrum vytváří nadprůměrné odměny ve srovnání s~celonárodním průměrem, avšak na mezinárodní úrovni zůstávají pod standardem Rakouska, Nizozemska nebo Lucemburska.
Finanční sektor je v~ČR charakteristický oligopolní strukturou se silným podílem zahraničních bank --- zaměstnavatelé v~tomto sektoru mají o~poznání vyšší ziskovost než v~jiných odvětvích, přesto \ac{KS} pokrytí je ve finančním sektoru velmi nízké.
Výsledkem je, že mzdová prémie finančního sektoru nad celostátním průměrem je v~\acloc{ČR} nižší než v~Rakousku nebo Švýcarsku, přestože produktivita na zaměstnance v~bankovnictví je komparabilní nebo vyšší --- přebytek jde primárně do dividendových výnosů zahraničních vlastníků.
Odvětvová \ac{KS} ve finančním sektoru by mohla zachytit část ziskového přebytku ve prospěch zaměstnanců --- nebyl by to administrativní zásah do trhu, ale institucionalizace kolektivní vyjednávací pozice zaměstnanců vůči oligopolně silným zaměstnavatelům.

Mapa průměrných mezd ve finančním sektoru a~pojišťovnictví (\ac{NACE} K) v~\acs{geo-EU27} umísťuje \acacc{ČR} do středního pole: Pražské finanční centrum vytváří nadprůměrné odměny ve srovnání s~celonárodním průměrem, avšak na mezinárodní úrovni zůstávají pod standardem Rakouska, Nizozemska nebo Lucemburska.
Finanční sektor je v~ČR charakteristický oligopolní strukturou se silným podílem zahraničních bank --- zaměstnavatelé v~tomto sektoru mají o~poznání vyšší ziskovost než v~jiných odvětvích, přesto \ac{KS} pokrytí je ve finančním sektoru velmi nízké.
Výsledkem je, že mzdová prémie finančního sektoru nad celostátním průměrem je v~\acloc{ČR} nižší než v~Rakousku nebo Švýcarsku, přestože produktivita na zaměstnance v~bankovnictví je komparabilní nebo vyšší --- přebytek jde primárně do dividendových výnosů zahraničních vlastníků.
Odvětvová \ac{KS} ve finančním sektoru by mohla zachytit část ziskového přebytku ve prospěch zaměstnanců --- nebyl by to administrativní zásah do trhu, ale institucionalizace kolektivní vyjednávací pozice zaměstnanců vůči oligopolně silným zaměstnavatelům.

Mapa průměrných mezd ve finančním sektoru a~pojišťovnictví (\ac{NACE} K) v~\acs{geo-EU27} umísťuje \acacc{ČR} do středního pole: Pražské finanční centrum vytváří nadprůměrné odměny ve srovnání s~celonárodním průměrem, avšak na mezinárodní úrovni zůstávají pod standardem Rakouska, Nizozemska nebo Lucemburska.
Finanční sektor je v~ČR charakteristický oligopolní strukturou se silným podílem zahraničních bank --- zaměstnavatelé v~tomto sektoru mají o~poznání vyšší ziskovost než v~jiných odvětvích, přesto \ac{KS} pokrytí je ve finančním sektoru velmi nízké.
Výsledkem je, že mzdová prémie finančního sektoru nad celostátním průměrem je v~\acloc{ČR} nižší než v~Rakousku nebo Švýcarsku, přestože produktivita na zaměstnance v~bankovnictví je komparabilní nebo vyšší --- přebytek jde primárně do dividendových výnosů zahraničních vlastníků.
Odvětvová \ac{KS} ve finančním sektoru by mohla zachytit část ziskového přebytku ve prospěch zaměstnanců --- nebyl by to administrativní zásah do trhu, ale institucionalizace kolektivní vyjednávací pozice zaměstnanců vůči oligopolně silným zaměstnavatelům.

\input{texparts/python/eu_odvetvove_mzdy_mapa_K}



